\begin{tikzpicture}[
  every node/.style={font=\footnotesize\sffamily},
  participant/.style={draw=TFGBlue, fill=TFGLightBlue, rounded corners=2pt,
    minimum width=2.15cm, minimum height=0.65cm, align=center},
  lifeline/.style={dashed, draw=TFGMuted},
  message/.style={-{Latex[length=2mm]}, semithick, draw=TFGAccent},
  return/.style={-{Latex[length=2mm]}, dashed, draw=TFGMuted},
  message label/.style={fill=white, inner xsep=2pt, inner ysep=1pt,
    text=black, align=center}
]
  \node[participant] (user) at (0,0) {Miembro};
  \node[participant] (web) at (2.8,0) {Aplicación web};
  \node[participant] (qualification) at (5.6,0) {Cualificación};
  \node[participant] (storage) at (8.4,0) {Documentos};
  \node[participant] (ai) at (11.2,0) {Proveedor\\de IA};
  \node[participant] (db) at (14,0) {Persistencia};
  \foreach \x in {0,2.8,5.6,8.4,11.2,14} \draw[lifeline] (\x,-0.35) -- (\x,-7.05);
  \draw[message] (0,-0.7) -- node[message label, above]{solicita una generación} (2.8,-0.7);
  \draw[message] (2.8,-1.35) -- node[message label, above]{inicia la generación del elemento} (5.6,-1.35);
  \draw[message] (5.6,-2) -- node[message label, above]{solicita los PDF asociados} (8.4,-2);
  \draw[return] (8.4,-2.65) -- node[message label, above]{documentos y metadatos} (5.6,-2.65);
  \draw[message] (5.6,-3.3) -- node[message label, above]{instrucción, esquema y documentos} (11.2,-3.3);
  \draw[return] (11.2,-3.95) -- node[message label, above]{salida estructurada o error} (5.6,-3.95);
  \draw[message] (5.6,-4.6) -- node[message label, above]{registra resultado y estado} (14,-4.6);
  \draw[return] (14,-5.25) -- node[message label, above]{registro confirmado} (5.6,-5.25);
  \draw[return] (5.6,-5.9) -- node[message label, above]{resultado o error revisable} (2.8,-5.9);
  \draw[return] (2.8,-6.55) -- node[message label, above]{muestra el estado de la generación} (0,-6.55);
\end{tikzpicture}
